\chapter{전원과 안전}\label{ch:44-power-safety}

지금까지는 신호와 논리였다. 이 장은 기계에 전기를 대는 쪽과, 사람이 다치지 않고 기계가 틀린 답을 내지 않게 하는
장치를 다룬다. 레일을 왜 셋으로 나누는지, 전원 장치의 크기를 어떻게 정하는지, 퓨즈를 왜 두 겹으로 두는지, 0~V를
왜 땅에 꼭 한 곳에서만 묶는지, 비상 정지와 인터록이 어떻게 생겼고 왜 프로그램이 아니라 접점으로 하는지가 차례로
나온다.

출처는 \code{hardware/common/power-unit.md}(전원부), 결정 0003의 6--7절(전원 레일, 접지와 보호),
\code{hardware/common/rotor-cartridge.md}(로터 교체와 인터록), 설계 개요(\code{design-brief.md})의 6절(안전)이다. 전원부는
아직 \emph{개념 설계} 단계다. 회로도도 시뮬레이션도 없고, 품번은 종류를 보이는 예이며, 많은 값이 추정이다. 이 장의
숫자도 문서의 표기를 따라 (계산), (추정), (미확인)을 붙인다. 앞의 두 장에서 본 보드처럼 ``시뮬레이션으로 확인했다''고 말할
수 있는 것은 이 장에 거의 없다.

\section{전원부 한눈에}\label{sec:p-overview}

\Why 건물의 콘센트는 220~V \term{교류}{AC}다. 전압의 방향이 1초에 60번씩 앞뒤로 오가고, 맨손으로 닿으면 위험하다. 이 기계의
회로는 24~V \term{직류}{DC}로 돈다. 그래서 전원부의 첫째 일은 위험한 전압을 \emph{인증된 상자 하나 안에 가두는 것}이다.
220~V는 접지 달린 인입구로 들어와 주 스위치(2극)와 차단기(2극 10~A)를 지나 곧바로 \term{전원 장치}{power supply unit,
PSU}로 들어가고, 거기서 나오는 것은 모두 극저전압이다. 전원 장치는 상용 전원과 출력 사이가 강화 절연된 스위칭 전원
장치이고, 출력이 SELV로 표시된 것을 고른다(규격 원문은 미확인). 관람객이 닿을 수 있는 것(드럼, 패널, 케이블)은 모두
24~V다. 극저전압의 상한은 교류 50~V, 직류 120~V다~\cite{iec61140}. 참고로 1940년대의 실물 봄베는 220~V 직류에 약
10~A를 썼다~\cite{us6812report1944}.

\General 노트북 충전기가 같은 일을 한다. 220~V는 충전기 상자 안에서만 다루고, 노트북으로 가는 선에는 20~V 안팎의 직류만
흐른다. 그래서 노트북의 금속 몸체를 만져도 위험하지 않다.

\Ours 전원부는 \emph{기계마다 하나}다(기어식에 하나, 릴레이식에 하나). 그 기계에 붙은 공통부(테스트 레지스터, 패널,
표시)도 같은 전원부에서 먹는다. 나가는 레일은 셋이다(표~\ref{tab:p-rails}, 그림~\ref{fig:p-unit}).

\begin{table}[ht]
\centering
\caption{세 레일(결정 0003 6절, 전원부 문서 1절).}
\label{tab:p-rails}
\small
\begin{tabularx}{\linewidth}{>{\raggedright\arraybackslash}p{22mm}>{\raggedright\arraybackslash}X>{\raggedright\arraybackslash}p{30mm}>{\raggedright\arraybackslash}p{28mm}}
\toprule
레일 & 먹이는 것 & 끊는 것 & 전원 장치 \\
\midrule
\code{+24V\_L} (논리, 상시) & 테스트 레지스터(유지·판정·램프), 표시 릴레이, 위치 표시, 감시 장치의 현장 쪽, 센서, 안전 고리, 팬, 조작부 & 주 스위치만 & PSU-L 24~V 10~A (두 기계 같음) \\
\code{+24V\_SW} (비상 정지로 끊김) & 뱅크의 타이밍 출력, 릴레이식의 계수기·시프터 코일과 순서기, 기어식 모터 드라이버의 허가 입력 & 비상 정지(K$_\mathrm{E}$) & 릴레이식: PSU-SW 24~V 27~A 이상, 6유닛마다 하나. 기어식: \code{+24V\_L}에서 갈라 냄(1~A 이하) \\
\code{+48V\_M} (기어식 모터) & 스테퍼 모터 드라이버. 구동부의 덮개 안에만 & 비상 정지(K$_\mathrm{E}$)와 덮개(K$_\mathrm{I}$) & PSU-M 48~V 7.3~A 이상 \\
\bottomrule
\end{tabularx}
\end{table}

\begin{figure}[tbp]
\centering
\resizebox{\linewidth}{!}{% 전원부의 레일과 안전 회로(한 줄 그림). 그림 fig:hw-power-unit
% 근거: hardware/common/power-unit.md 3-4절, docs/decisions/0003-interfaces.md 4.6, 6, 7절.
\begin{tikzpicture}[x=1cm,y=1cm,
    psu/.style={draw,rounded corners=1.5pt,minimum width=24mm,minimum height=11mm,align=center,inner sep=2pt,font=\scriptsize},
    ld/.style={draw=cgray,fill=white,rounded corners=1.5pt,align=left,inner sep=2.5pt,font=\tiny},
    km/.style={draw=cred,fill=credbg,minimum width=13mm,minimum height=6mm,inner sep=1pt,font=\tiny,align=center},
    sr/.style={draw=cred,fill=white,minimum width=11mm,minimum height=7mm,inner sep=1pt,font=\scriptsize,align=center},
    rail/.style={line width=1.4pt},
    lbl/.style={font=\tiny,inner sep=1pt}]
  % 상용 전원
  \node[psu,fill=clight,draw=cgray,minimum height=16mm] (ac) at (0,0) {상용 전원\\인입구(접지)\\주 스위치\\차단기 10 A};
  % 전원 장치
  \node[psu,fill=cbluebg,draw=cblue] (pl) at (3.0,2.6) {PSU-L\\24.5 V 10 A};
  \node[psu,fill=corangebg,draw=corange] (ps) at (3.0,0) {PSU-SW\\24.5 V 27 A\\{\tiny 릴레이식, 6유닛마다}};
  \node[psu,fill=credbg,draw=cred] (pm) at (3.0,-2.6) {PSU-M\\48 V 7.3 A\\{\tiny 기어식}};
  \draw[wire] (ac.east) -- ++(0.45,0) coordinate (t);
  \draw[wire] (t) |- (pl.west);
  \draw[wire] (t) -- (ps.west);
  \draw[wire] (t) |- (pm.west);
  % +24V_L
  \draw[rail,cblue] (pl.east) -- (13.4,2.6);
  \node[lbl,text=cblue,anchor=south west] at (4.3,2.65) {\code{+24V\_L}};
  \node[ld,anchor=south] (l1) at (11.6,3.25) {테스트 레지스터, 표시 릴레이\\감시 장치, 센서, 팬, 조작부\\F 축 브레이크(기어식, 덮개 스위치 경유)};
  \draw[wire] (11.6,2.6) -- (l1.south);
  \node[lbl,anchor=west] at (11.65,2.93) {갈래마다 퓨즈};
  % 안전 릴레이
  \node[sr] (ke) at (7.3,1.3) {K$_\mathrm{E}$};
  \node[sr] (ki) at (8.7,1.3) {K$_\mathrm{I}$};
  \draw[wire] (7.3,2.6) -- (ke.north);
  \draw[wire] (8.7,2.6) -- (ki.north);
  \node[ld,anchor=south] (lp1) at (7.0,3.25) {비상 정지 고리\\제어 커넥터 22--23};
  \node[ld,anchor=south] (lp2) at (9.0,3.25) {인터록 고리\\24--25};
  \draw[wire,dashed] (7.3,2.6) -- (7.3,3.25);
  \draw[wire,dashed] (8.7,2.6) -- (8.7,3.25);
  % +24V_SW
  \node[km] (k1) at (5.6,0) {KM1·KM2};
  \draw[rail,corange] (ps.east) -- (k1.west);
  \draw[rail,corange] (k1.east) -- (9.2,0);
  \node[lbl,text=corange,anchor=south west] at (6.3,0.05) {\code{+24V\_SW}};
  \node[ld,anchor=west] (s1) at (9.4,0.32) {퓨즈 5 A × 6 → 계수기 유닛 보드 → 시프터 4장};
  \node[ld,anchor=west] (s2) at (9.4,-0.32) {퓨즈 → 공통 계수기; 퓨즈 → K$_\mathrm{I}$ 접점 → 클럭 순서기};
  \draw[wire] (9.2,0) |- (s1.west);
  \draw[wire] (9.2,0) |- (s2.west);
  % +48V_M
  \node[km] (k2) at (5.6,-2.6) {KM1·KM2};
  \node[km] (k3) at (7.7,-2.6) {KM3};
  \draw[rail,cred] (pm.east) -- (k2.west);
  \draw[rail,cred] (k2.east) -- (k3.west);
  \draw[rail,cred] (k3.east) -- (9.8,-2.6);
  \node[ld,anchor=west] (m1) at (9.8,-2.6) {스테퍼 드라이버 (구동부의 덮개 안)};
  \node[lbl,text=cred,anchor=south west] at (8.6,-2.55) {\code{+48V\_M}};
  \node[lbl,anchor=north west,align=left,text=cgray] at (4.5,-3.05) {기어식의 \code{+24V\_SW}(1 A 이하)는 \code{+24V\_L}에서\\KM1·KM2의 다른 극을 거쳐 갈라 낸다};
  % 안전 릴레이가 접촉기를 구동
  \draw[->,cred,dashed,line width=0.6pt] (ke.south) -- (7.3,-1.3) -- (5.6,-1.3) -- (k2.north);
  \draw[->,cred,dashed,line width=0.6pt] (5.6,-1.3) -- (k1.south);
  \draw[->,cred,dashed,line width=0.6pt] (ki.south) -- (8.7,-1.7) -- (7.7,-1.7) -- (k3.north);
  \node[lbl,text=cred,anchor=north west] at (5.7,-1.36) {비상 정지: 접촉기 둘을 직렬로};
  \node[lbl,text=cred,anchor=west,align=left] at (8.8,-1.5) {덮개: 모터만 (릴레이식: 카트리지 고리, 순서기만)};
  % 0 V와 PE
  \draw[rail,cgray] (1.6,-4.4) -- (7.0,-4.4);
  \node[lbl,anchor=south west] at (1.6,-4.35) {0 V 모선: 모든 전원 장치의 0 V};
  \draw[rail,cgreen] (8.6,-4.4) -- (13.4,-4.4);
  \node[lbl,anchor=south west,text=cgreen] at (8.6,-4.35) {PE 모선: 프레임, 덮개, 케이스, 모터};
  \draw[line width=1pt,double] (7.0,-4.4) -- (8.6,-4.4);
  \node[lbl,anchor=north] at (7.8,-4.5) {연결 고리 (한 곳뿐)};
  \draw[wire,cgreen] (ac.south) |- (1.0,-4.9) -| (13.2,-4.4);
  \node[lbl,text=cgreen,anchor=south west] at (1.0,-4.88) {PE};
\end{tikzpicture}
}
\caption{전원부의 레일과 안전 회로(보고서와 같은 그림, 전원부 문서 3--4절). 굵은 선이 레일이고, 빨간 점선은 안전 릴레이가
접촉기를 움직이는 길이다. 아래의 0~V 모선과 PE 모선은 연결 고리 하나로만 이어진다(\ref{sec:p-earth}절).}
\label{fig:p-unit}
\end{figure}

\section{왜 레일을 나누는가}\label{sec:p-rails}

\Why 레일을 나누는 이유는 하나다. \emph{끊어야 하는 때가 서로 다르다}. 비상 정지 버튼을 누르면 움직임과 판정은 멈춰야
하지만, 판정의 \emph{결과}는 남아 있어야 한다. 정지가 났는데 비상 정지와 함께 램프가 다 꺼지면 무엇이 일어났는지 알 수
없다. 그래서 움직임과 박자를 만드는 쪽(\code{+24V\_SW}, \code{+48V\_M})은 끊고, 결과를 붙들고 보이는 쪽(\code{+24V\_L})은
살려 둔다. 타이밍 신호가 사라지면 테스트 레지스터의 \code{STROBE} 접점이 열리지 않으므로 시험 전압도 걸리지 않는다.
비상 정지가 시험 전압을 따로 끊을 필요가 없는 이유다.

기어식의 덮개는 한 단계 더 좁다. 덮개를 열면 모터의 48~V만 끊고 시험 전압과 레지스터는 살려 둔다. 덮개를 연 채 손
크랭크로 드럼을 돌려 정지가 나는 것을 보이는 시연이 그대로 되게 하려는 것이다.

\General 엘리베이터의 비상 정지는 모터를 세우지만 칸 안의 조명과 비상 통화 장치는 그대로 둔다. 집의 분전반이 조명
회로와 콘센트 회로를 따로 두는 것도, 한쪽에 문제가 생겨도 다른 쪽은 살아 있게 하려는 것이다.

\Ours 이 나눔에는 대가도 있다. 감지 릴레이 $K$는 \code{+24V\_L}에 있어서 비상 정지에도 자기 유지를 놓지 않는다. 그래서
위치 주기 한가운데서 비상 정지가 나면 지워지지 않은 $K$가 남고, 결정 0003은 ``\code{HOME}은 끝에 \code{CLEAR}를 정확히
한 번 낸다''는 규칙을 더했다(\ref{ch:boardverify}장의 실제 사례). 그 밖의 원칙은 이렇다.
\begin{itemize}
\item \textbf{48~V는 덮개 안에만.} \code{+48V\_M}은 구동부의 덮개 안에만 있고 어떤 인터페이스 커넥터에도 나오지 않는다.
단자는 24~V와 꽂히지 않는 형식(7.62~mm 3극: +48~V, 0~V$_\mathrm{M}$, PE)이다. 48~V도 직류 60~V 이하의 범위로 보지만 규격
원문은 미확인이다.
\item \textbf{전원 장치를 나란히 잇지 않는다.} 릴레이식 12유닛에서는 PSU-SW가 둘인데, 두 출력을 한 모선에 묶지 않는다.
각자 자기 여섯 유닛을 먹이고 0~V만 한 점에서 만난다.
\item \textbf{기어식의 \code{+24V\_SW}.} 슬롯 센서의 출력단과 드라이버 허가 입력뿐이어서(0.2~A 이하, 추정) 따로 전원
장치를 두지 않고 \code{+24V\_L} 모선에서 접촉기의 한 극씩을 거쳐 갈라 낸다. 슬롯 센서 자체는 \code{+24V\_L}에서 먹고
출력단(TBD62783A)만 \code{+24V\_SW}에서 먹으므로, 비상 정지가 타이밍 출력만 정확히 없앤다.
\end{itemize}

\section{전류 예산: 전원 장치의 크기}\label{sec:p-budget}

\Why 전원 장치는 정해진 전류까지만 낸다. 부하가 그보다 많이 끌면 장치는 전압을 낮춰 전류를 제한하거나 아예 꺼진다.
릴레이 코일은 전압이 떨어지면 붙지 않거나 놓이므로, 전원이 모자란 릴레이 기계는 ``가끔 틀리는 기계''가 된다. 그래서
부하를 하나씩 더해 \emph{최악}을 구하고, 그보다 여유 있는 장치를 고른다. 평균이 아니라 최악인 이유는 최악이 실제로
일어나기 때문이다.

\General 멀티탭에 전열기와 헤어드라이어를 함께 꽂으면 각각은 괜찮아도 합이 멀티탭의 정격을 넘는다. 늘 동시에 쓰지
않더라도 동시에 쓸 수 있으면 그 합으로 정격을 정해야 한다.

\Ours 레일마다 예산은 이렇다(전원부 문서 2절. 코일은 200~mW로 보았고 미확인이다).
\begin{itemize}
\item \textbf{\code{+24V\_L}.} 테스트 레지스터 0.47~A, 표시 릴레이 1.6~A(한도), 위치 표시 보드 0.33~A, 팬 0.8~A, 조작부
0.8~A 등을 더해 최대 기어식 5.2~A, 릴레이식 4.8~A다(일부는 추정). PSU-L 10~A(240~W)이면 정격의 약 50\%다.
\item \textbf{\code{+24V\_SW}(릴레이식).} 유닛 하나가 최악 3.7~A이고(\ref{sec:el-power}절), 유닛 여섯 22.2~A에 공통 계수기
0.25~A와 순서기 0.5~A(추정)를 더해 23.0~A다. PSU-SW는 27~A 이상(650~W급)을 6유닛마다 하나 둔다. 최악이 정격의 85\%,
평균이 35--49\%다.
\item \textbf{\code{+48V\_M}(기어식).} 모터의 기계적 출력은 연속 45~W, 두 자리올림이 겹치는 순간 106~W다. 드라이버와
모터의 손실을 더해 전원에서 연속 2--4~A, 순간 6~A 안팎을 끈다고 보고(추정) PSU-M을 7.3~A 이상(350~W급)으로 잡았다.
\end{itemize}

몇 가지 결정은 이 예산에서 나온다.
\begin{itemize}
\item \textbf{평균으로 잡지 않은 이유.} 평균으로 잡으면 PSU-SW가 480~W(20~A)로 된다. 그러나 오프셋 다이얼을 모두 같은
값(15나 23)으로 두면 모든 유닛의 시프터 네 장이 동시에 넷 단씩 켜진다. 평균으로 잡은 기계는 다이얼 설정에 따라 전압이
꺼지는 기계가 된다. 하나로 합쳐 960~W(40~A)를 쓰면 12유닛의 최악 45.2~A에 모자라므로 6유닛마다 하나씩 둔다. 그러면
6유닛 단계에서는 하나만 사면 된다.
\item \textbf{모르는 값이 크기를 바꾼다.} 코일이 200~mW가 아니라 140~mW이면 최악이 약 17~A로 줄어 PSU-SW가 480~W급으로
작아질 수 있다. 시제품에서 코일 전류를 재고 다시 정한다.
\item \textbf{48~V의 커패시터는 접촉기 앞에.} PSU-M의 출력에 4,700~$\mu$F 커패시터를 두는데, 접촉기의 \emph{전원 쪽}에
둔다. 접촉기 뒤에 두면 접촉기가 열린 뒤에도 그 커패시터가 모터를 계속 먹인다. 앞에 두면 접촉기가 열린 순간 드라이버에
남는 것은 드라이버 안의 작은 커패시터뿐이어서, 덮개를 연 뒤 0.1초 안에 선다는 기준을 지킨다. 모터와 드럼이 되돌려 주는
에너지는 0.2~J 이하(계산)여서 제동 저항을 두지 않는다. 전원 장치는 과전류에서 꺼졌다 켜지기를 되풀이하는 형이 아니라
전류를 제한하는 형을 고른다(스테퍼의 순간 전류, 추정).
\item \textbf{상용 전원의 소비와 열.} 릴레이식은 평균 약 0.4~kW(6유닛)와 0.7~kW(12유닛), 최대 약 0.8과 1.4~kW, 기어식은 최대
약 0.5~kW다. 전원 장치 자신도 약 10\%를 열로 버린다. 릴레이식 12유닛에서 평균 약 70~W이고, 패널 칸 뒤의 흡기로 식힌다.
\end{itemize}

\paragraph{전압의 예산.} 전류만이 아니라 전압도 예산이 있다. 가장 무거운 갈래(릴레이식 유닛 하나, 3.7~A, 편도 2.5~m)에서
전원 장치부터 보드 단자까지 잃는 전압을 더하면 표~\ref{tab:p-drop}와 같다. 전원 장치를 24.5~V $\pm$ 0.5로 맞추면 보드
단자에서 23.6--24.6~V이고, 요구는 23.0--26.4~V다(\ref{sec:el-series}절). 접촉기와 퓨즈의 저항은 추정이고 시제품에서 잰다.

\begin{table}[ht]
\centering
\caption{가장 무거운 갈래의 전압 강하(전원부 문서 6절, 계산).}
\label{tab:p-drop}
\small
\begin{tabular}{lr}
\toprule
구간 & 강하 \\
\midrule
접촉기 두 극(극당 약 2~m$\Omega$, 23~A) & 0.09 V \\
퓨즈 5~A(약 12~m$\Omega$)와 단자 & 0.06 V \\
전선 1.5~mm$^2$ 왕복 5~m & 0.22 V \\
\midrule
보드 단자까지 & \textbf{0.37 V} \\
\bottomrule
\end{tabular}
\end{table}

\section{퓨즈와 자동 복귀 퓨즈}\label{sec:p-fuse}

\Why 단락이 나면 전류가 전원이 낼 수 있는 만큼 커진다(\ref{sec:el-vir}절). 전선은 전류에 비례해 달아오르므로, 그 전류를
누군가 끊어야 한다. 끊는 부품은 두 종류다.
\begin{itemize}
\item \term{퓨즈}{fuse}는 가는 금속선이다. 정해진 전류를 넘는 전류가 일정 시간 흐르면 녹아 끊긴다. 확실하게, 영구히 끊고,
갈아 끼워야 다시 쓴다. \term{시간 지연형}{time-delay}은 전원을 켤 때의 짧은 큰 전류(돌입 전류)는 견디고 오래 계속되는
과전류에만 끊긴다.
\item \term{자동 복귀 퓨즈}{PTC resettable fuse}는 열을 받으면 저항이 급격히 커지는 고분자 부품이다. 과전류로 달아오르면
저항이 뛰어 전류가 작은 값으로 줄고, 원인이 사라져 식으면 저절로 원래대로 돌아온다. 갈아 끼울 필요가 없는 대신, 끊긴
동안에도 작은 전류가 흐르고(완전히 끊지는 않는다) 퓨즈보다 느리다.
\end{itemize}
둘은 지키는 대상이 다르다. 그래서 두 겹으로 둔다(그림~\ref{fig:p-fuse}). \emph{보드 안의 고장은 보드의 자동 복귀 퓨즈가,
전선의 고장은 전원부의 퓨즈가} 맡는다. 보드 안에서 부품이 단락되면 보드 입구의 자동 복귀 퓨즈가 전류를 줄이고, 고장이
사라지면 저절로 돌아온다. 전원부에서 보드까지의 전선이 틀에 닿아 단락되면 그 단락은 보드 퓨즈보다 앞에 있으므로 보드
퓨즈는 아무것도 할 수 없다. 이것은 전원부 단자대의 퓨즈가 끊는다. 그러려면 전원부 퓨즈는 보드의 자동 복귀 퓨즈보다
\emph{크고}(보드 고장에 먼저 끊기지 않게) 전선이 견디는 전류보다 \emph{작아야}(전선이 타기 전에 끊기게) 한다.

\begin{figure}[ht]
\centering
\begin{tikzpicture}[x=1cm,y=1cm]
  \node[blkb,minimum width=18mm] (psu) at (0,0) {전원 장치\\{\scriptsize 전류 제한}};
  \node[blkr,minimum width=18mm] (f) at (2.9,0) {갈래 퓨즈\\{\scriptsize 단자대, 끊김 표시}};
  \draw[line width=1.6pt,cgray] (f.east) -- node[above,font=\scriptsize] {전선 0.5--1.5 mm$^2$} (7.0,0);
  \node[blko,minimum width=18mm] (ptc) at (8.1,0) {자동 복귀\\퓨즈};
  \node[blkn,minimum width=26mm] (brd) at (10.6,0) {보드의 회로\\{\scriptsize 역극성 보호, TVS, …}};
  \draw (psu) -- (f); \draw (ptc) -- (brd);
  \draw[cgray,dashed,rounded corners=2pt] (7.0,-0.75) rectangle (12.1,0.75);
  \node[note,anchor=south] at (9.55,0.75) {보드};
  \draw[cgray,dashed,rounded corners=2pt] (-1.1,-0.75) rectangle (4.2,0.75);
  \node[note,anchor=south] at (1.55,0.75) {전원부};
  % faults
  \draw[cred,line width=1pt,dashed] (5.6,0) -- (5.6,-1.3);
  \node[note,text=cred,anchor=north,align=center] at (5.6,-1.3) {전선이 틀에 닿음\\$\to$ 갈래 퓨즈가 끊는다};
  \draw[cred,line width=1pt,dashed] (10.6,-0.45) -- (10.6,-1.3);
  \node[note,text=cred,anchor=north,align=center] at (10.6,-1.3) {보드 안의 단락\\$\to$ 자동 복귀 퓨즈가 줄인다};
  \node[note,anchor=north,align=center] at (2.9,-1.3) {보드 것보다 크고\\전선이 견디는 것보다 작게};
\end{tikzpicture}
\caption{보호의 두 겹. 고장이 어디서 났느냐에 따라 맡는 부품이 다르다.}
\label{fig:p-fuse}
\end{figure}

\General 집의 분전반에는 회로마다 차단기가 있고, 가전제품 안에는 따로 퓨즈가 있다. 차단기는 벽 속의 전선을, 가전의 퓨즈는
가전을 지킨다. 헤어드라이어가 고장 나면 드라이어의 퓨즈가 끊기고 집의 다른 콘센트는 그대로다. 컴퓨터의 USB 포트에는
자동 복귀 퓨즈가 있어서, 고장 난 기기를 꽂았다 빼면 포트가 저절로 살아난다.

\Ours 전원부 문서 6절의 표에서 몇 줄을 옮기면 이렇다.

\begin{table}[ht]
\centering
\caption{갈래 퓨즈의 예(전원부 문서 6절). 모두 시간 지연형이다.}
\label{tab:p-fuses}
\small
\begin{tabularx}{\linewidth}{l>{\raggedright\arraybackslash}X r >{\raggedright\arraybackslash}p{30mm} r}
\toprule
퓨즈 & 갈래 & 정격 & 보드 안의 보호 & 전선 \\
\midrule
F11 & 테스트 레지스터 & 4 A & 논리 1.1 A + 표시 2 A & 0.75 mm$^2$ \\
F12 & 위치 표시 보드 (릴레이식) & 2 A & 0.5 A (최대 부하 0.33 A) & 0.75 mm$^2$ \\
F18, F19 & 비상 정지 고리, 인터록 고리 & 0.5 A & --- & 0.5 mm$^2$ \\
F31--F36 & 계수기 유닛 보드 (유닛마다 하나) & 5 A & 4 A & 1.5 mm$^2$ \\
F38 & 클럭 순서기 (K$_\mathrm{I}$ 접점 경유) & 3.15 A & 3 A & 0.75 mm$^2$ \\
F51 & \code{+48V\_M} (기어식) & 10 A & 드라이버 & 1.5 mm$^2$ \\
\bottomrule
\end{tabularx}
\end{table}

\begin{itemize}
\item \textbf{보드의 자동 복귀 퓨즈}는 유지 전류(끊기지 않고 계속 흘릴 수 있는 전류)를 최대 부하의 1.3--2배로 잡는다(결정 0003
7.3절). 위치 표시 보드는 최대 0.33~A에 0.5~A다.
\item \textbf{유닛마다 퓨즈 하나.} 시프터 보드에는 전원 단자가 없고 코일 전류가 계수기 유닛 보드의 퓨즈 뒤에서 \code{CTRL}
선으로 온다. 그래서 유닛 하나의 배선에 고장이 나면 그 유닛의 퓨즈만 끊기고 나머지 유닛의 전원은 그대로다.
\item \textbf{느리게 끊기는 갈래.} 전원 장치의 전류 제한이 퓨즈를 끊을 만큼 크지 않은 갈래(PSU-L 10~A에 F11 4~A)는 끊기는
데 몇 초가 걸린다(추정). 전자식 차단기로 바꾸는 것을 구매 때 검토한다.
\item \textbf{시험 전압 선.} 테스트 레지스터의 시험 전압은 보드의 F1(1.1~A)을 지난다. 메뉴 배선망의 어느 선이 금속 틀이나
0~V에 닿으면 F1이 끊기고, 그 위치는 \emph{거짓 정지}가 된다. 틀리더라도 안전한 쪽으로 틀린다. 직렬 저항으로 전류를 묶지 않는
것은 저항이 감지 릴레이의 코일 전압을 깎기 때문이다(\ref{sec:el-series}절).
\end{itemize}

\section{보드 입구의 네 부품}\label{sec:p-input}

\Why 전원 단자가 있는 모든 보드는 입구에 같은 네 부품을 같은 순서로 둔다(결정 0003 7.3절, 그림~\ref{fig:p-input}).
\begin{enumerate}
\item \textbf{자동 복귀 퓨즈.} 보드 안의 고장을 맡는다(\ref{sec:p-fuse}절).
\item \textbf{역극성 보호: 직렬 P채널 MOSFET.} 전원선을 거꾸로 꽂아도 보드가 상하지 않게 한다. 가장 쉬운 방법은 다이오드
하나를 직렬로 넣는 것이지만, 그러면 보드 전체가 0.7~V를 잃는다(\ref{sec:el-diode}절). 23.0~V 하한에서 0.7~V는 크다. MOSFET은
켜지면 작은 저항일 뿐이어서 강하가 0.1~V 이하다.
\item \textbf{TVS 33~V.} $+24$~V와 0~V 사이에서, 전원선에 실려 온 순간의 튐을 33~V 근처에서 삼킨다. 정상 범위(최대 26.4~V)에서는
흐르지 않는다.
\item \textbf{전해 커패시터 100~$\mu$F 이상.} 저수지다(\ref{sec:el-cap}절). 릴레이가 한꺼번에 붙는 순간의 전류를 받아 준다.
\end{enumerate}

MOSFET이 방향을 가리는 원리는 이렇다. MOSFET 안에는 만들 때 생기는 다이오드(\term{몸체 다이오드}{body diode})가 드레인에서
소스 쪽으로 하나 들어 있다. 드레인을 전원 쪽에 두고 게이트를 100~k$\Omega$로 0~V에 묶는다.
\begin{itemize}
\item \textbf{바로 꽂으면} 먼저 몸체 다이오드로 전류가 흘러 소스가 약 23~V가 된다. 게이트는 0~V이므로 게이트가 소스보다 20~V
넘게 낮다. P채널이므로 확실히 켜지고, 그 뒤로는 전류가 다이오드가 아니라 켜진 채널로 흘러 0.7~V를 잃지 않는다. 게이트와
소스 사이의 제너 12~V가 $V_{GS}$를 $-12$~V에서 묶어 게이트의 한계를 지킨다.
\item \textbf{거꾸로 꽂으면} 드레인 쪽이 보드의 0~V 선보다 24~V 낮아져 몸체 다이오드가 막힌다. 게이트는 100~k$\Omega$로 보드의
0~V 선(이제 실제로는 높은 쪽)에 묶여 있어 소스보다 낮아지지 않으므로 채널도 꺼진 채다. 전류가 0이다.
\end{itemize}

\begin{figure}[ht]
\centering
\begin{circuitikz}[x=1cm,y=1cm]
  \ctikzset{bipoles/length=9mm}
  \draw (0,2.6) node[left,font=\footnotesize,align=right] {단자 1\\$+24$ V} to[thRp,l={자동 복귀 퓨즈}] (2.4,2.6);
  \node[pmos,rotate=-90,xscale=-1,anchor=D] (Q) at (2.4,2.6) {};
  \node[font=\scriptsize,anchor=south] at ([yshift=3mm]Q.center) {P채널};
  \draw (Q.S) -- (5.2,2.6) coordinate (s);
  \draw (Q.G) -- (Q.G |- 0,1.55) coordinate (gg) -- (4.6,1.55);
  \draw (4.6,1.55) to[zD*,l_={12 V}] (4.6,2.6);
  \node[circ] at (4.6,2.6) {};
  \draw (gg) to[R,l_={100 k$\Omega$},*-] (gg |- 0,0);
  \draw (s) -- (6.2,2.6) coordinate (t);
  \draw (6.2,0) to[tvsD*,l={TVS 33 V}] (t);
  \node[circ] at (t) {};
  \draw (t) -- (8.0,2.6) coordinate (c);
  \draw (c) to[C,l={100 $\mu$F},*-*] (8.0,0);
  \draw (c) -- (9.6,2.6) node[right,font=\footnotesize] {보드의 $+24$ V};
  \draw (0,0) node[left,font=\footnotesize,align=right] {단자 2\\0 V} -- (9.6,0) node[right,font=\footnotesize] {보드의 0 V};
  \node[note,anchor=north] at (2.4,-0.1) {보드 안의 고장};
  \node[note,anchor=north] at (3.6,-0.55) {거꾸로 꽂음};
  \node[note,anchor=north] at (6.2,-0.1) {순간의 튐};
  \node[note,anchor=north] at (8.0,-0.55) {순간의 큰 전류};
\end{circuitikz}
\caption{전원 단자가 있는 보드의 입구(결정 0003 7.3절). 왼쪽부터 자동 복귀 퓨즈, 직렬 P채널 MOSFET(드레인이 전원 쪽, 게이트에
제너 12~V와 100~k$\Omega$), TVS 33~V, 커패시터. 아래 글은 각 부품이 맡는 일이다.}
\label{fig:p-input}
\end{figure}

\General 건전지를 넣는 장난감 가운데 거꾸로 넣으면 아예 접점이 닿지 않게 만든 것이 있다. 기구로 막는 역극성 보호다. 전자
기기는 단자를 그렇게 만들 수 없을 때 이런 회로로 막는다.

\Ours 결정 0003이 나오기 전에 만든 테스트 레지스터와 단위 시험기에는 자동 복귀 퓨즈와 TVS만 있었다. 전원을 거꾸로 꽂으면
TVS가 순방향 다이오드로 도통해 자동 복귀 퓨즈가 끊길 때까지 그 전류를 받는다. 동작은 하지만 TVS가 그 전류를 얼마나 오래
견디는지 확인하지 못했다. 그래서 P채널 MOSFET을 더하기로 했다. 전원 단자의 핀도 ``1 = +24~V, 2 = 0~V''로 모든 보드에서
통일하고 기판의 실크에 적는다. 사람이 거꾸로 꽂을 일 자체를 줄이는 것이다. 시제품에서는 ``거꾸로 꽂았을 때 손상 없음, 전류
0''을 확인한다(결정 0003의 I8).

\section{인터페이스의 약속}\label{sec:p-rules}

\Why 기계는 보드 수십 장과 케이블 수백 가닥으로 이루어지고, 관람객이 직접 꽂는 케이블도 있다. 사고는 대개 부품 안이 아니라
\emph{이음매}에서 난다. 그래서 결정 0003은 커넥터와 케이블에 ``어떻게 꽂아도 위험한 일이 일어나지 않게'' 하는 약속을 두었다.
\begin{itemize}
\item \textbf{전원은 신호 커넥터로 보내지 않는다.} 보드마다 전원 단자(5.08~mm 2극)를 따로 둔다. D-sub 핀 하나는 1~A 이하,
커넥터 하나의 전원은 2~A 이하다. 릴레이식 유닛 하나는 최악 2.5--3.6~A를 쓰므로 신호 커넥터의 핀 몇 개로는 받을 수 없다.
\item \textbf{26선 커넥터에는 전원 핀이 없다.} 26선은 모두 신호선이다. 메뉴 케이블을 어디에 어떻게 꽂아도 전원끼리 닿지
않는다.
\item \textbf{보드는 수(핀), 케이블은 양 끝이 암(소켓).} 풀려 있는 케이블 끝에 드러난 핀이 없다. 한쪽이 꽂힌 채 다른 쪽이
금속 패널에 닿아도 선끼리 이어지지 않는다. 관람객이 다루는 것은 케이블이다.
\item \textbf{48~V는 인터페이스에 나오지 않는다}(\ref{sec:p-rails}절).
\item \textbf{역기전력은 받는 쪽이 자기 안에서 처리한다.} 코일을 다는 입력은 핀과 코일 사이에 직렬 다이오드를 두고, 코일의
클램프는 그 다이오드의 코일 쪽에 두며, 핀 쪽에는 풀다운을 둔다.
\end{itemize}

마지막 약속은 이유가 길다(\ref{sec:relay-flyback}절). 테스트 레지스터의 타이밍 릴레이는 빨리 놓이도록 다이오드 + 제너 30~V로
클램프한다. 처음 회로에서는 코일의 $+$ 쪽이 곧 커넥터 핀이어서, 신호가 꺼지는 순간 핀이 약 $-31$~V까지 내려갔다. 그 음전압이
케이블을 타고 내는 쪽으로 가면 두 가지가 깨진다. 기어식의 광센서 출력(내압 30~V급, 미확인)에는 $24 + 31 = 55$~V가 걸려 견디지
못한다. 출력에 클램프 다이오드가 든 고측 구동기(TBD62783A 같은 것)는 그 음전압을 0.7~V에서 잘라 버려 제너의 효과를 없앤다.
그러면 릴레이의 복귀가 3--10~ms로 늘어 \ref{ch:relay-timing}장의 시간 가정이 깨진다. 핀과 코일 사이의 직렬 다이오드는 꺼지는
순간 막혀서, 코일의 전류가 보드 안의 클램프로만 돌게 하고 핀은 풀다운으로 0~V에 머물게 한다. 대가는 코일 전압 0.7~V다.

\General 노트북 충전기의 플러그와 휴대폰 충전기의 플러그는 모양이 달라 서로 꽂히지 않는다. 전압이 다른 것은 아예 꽂히지
않게 만드는 것이 가장 확실한 약속이다. 48~V 단자를 24~V와 다른 형식으로 고른 것이 같은 생각이다.

\Ours 이 약속들은 문서에만 있지 않다. 음성 시험(\ref{sec:v-mutants}절)이 ``전원 입력의 보호 소자 방향을 뒤집음'', ``신호
커넥터의 예약 핀에 +24~V'', ``장착 구멍을 0~V에'' 같은 위반을 일부러 만들어 검증기가 잡는지 본다.

\section{접지: 0 V, PE, 섀시}\label{sec:p-earth}

\Why 접지라는 말에는 서로 다른 세 넷이 섞여 있다. 하나씩 떼어 보자.
\begin{itemize}
\item \textbf{0~V}는 회로의 기준이자 전류의 복귀선이다(\ref{sec:el-vir}절). 회로도의 GND가 이것이다.
\item \term{보호접지}{protective earth, PE}는 건물의 땅에 이어진 선이다. 상용 전원 플러그의 접지 핀으로 들어온다. 금속 틀,
덮개, 전원 장치의 케이스, 모터 몸통을 여기에 잇는다. 220~V 선이 상자 안에서 금속 케이스에 닿으면 큰 전류가 PE로 흘러
차단기나 건물의 누전 차단기가 떨어진다. 케이스가 220~V가 된 채 사람을 기다리지 않게 하는 것이다.
\item \textbf{섀시}(넷 \code{CHASSIS})는 보드의 D-sub 금속 껍데기와 장착 구멍이다. 금속 기둥으로 틀(PE)에 닿는다. 0~V가 아니다.
\end{itemize}

이제 두 질문이 남는다. 0~V를 PE에 묶어야 하는가, 그리고 몇 곳에서 묶는가.

\paragraph{묶는다.} 0~V를 땅에서 띄워 두는 방식(SELV)도 있다. 이 프로젝트는 0~V를 PE에 묶는다(PELV). 이유는 고장이 \emph{드러나게}
하기 위해서다(그림~\ref{fig:p-fault}). 모든 코일의 한쪽이 0~V이고 스위치는 $+$ 쪽에 있다. 0~V가 땅에 묶여 있으면 $+$ 쪽 선이
틀에 닿는 첫 고장이 곧 단락이어서 퓨즈를 끊는다. 고장이 그 자리에서 드러난다. 0~V가 떠 있으면 첫 고장은 아무 일도 일으키지
않아 아무도 모른다. 그 상태로 둘째 고장이 나면 두 고장이 틀을 거쳐 접점을 건너뛰는 길을 만들고, 릴레이가 \emph{저절로} 붙을
수 있다. 정지 릴레이가 저절로 붙으면 거짓 정지이고, 통과 릴레이가 저절로 붙으면 \emph{놓친 정지}다. 기계 제어 회로의 일반
원칙이기도 하다(IEC 60204-1의 제어 회로 접지, 원문 미확인).

\begin{figure}[ht]
\centering
\begin{circuitikz}[x=1cm,y=1cm]
  \ctikzset{bipoles/length=8mm}
  \foreach \dx/\ttl in {0/{(가) 0 V를 한 점에서 PE에 묶음},7.6/{(나) 0 V가 떠 있음}} {
  \begin{scope}[xshift=\dx cm]
    \node[font=\small\bfseries] at (2.6,3.9) {\ttl};
    \draw (0,3) node[left,font=\footnotesize] {$+24$ V} to[fuse,l={F}] (1.6,3) -- (2.0,3) to[nos,l={접점}] (3.4,3) -- (4.2,3);
    \draw (4.2,3) to[L,l={코일}] (4.2,1.2) -- (0,1.2) node[left,font=\footnotesize] {0 V};
    \node[circ] at (4.2,3) {};
    \fill[cgray!30] (-0.3,-0.05) rectangle (5.5,0.2);
    \node[font=\scriptsize,anchor=north] at (2.35,-0.05) {금속 틀 (PE)};
  \end{scope}}
  % (가) XPE link and one fault
  \draw[line width=1.2pt,cgreen] (0.6,1.2) -- (0.6,0.2);
  \node[font=\scriptsize,text=cgreen,anchor=west] at (0.65,0.7) {연결 고리};
  \draw[cred,line width=1.2pt,dashed] (1.8,3) -- (1.8,0.2);
  \node[font=\scriptsize,text=cred,anchor=west,align=left] at (1.9,2.2) {첫 고장:\\$+$ 선이 틀에 닿음};
  \node[note,text=cred,anchor=north,align=center] at (2.35,-0.45) {$+24$ V $\to$ F $\to$ 틀 $\to$ 0 V: 단락\\퓨즈가 끊겨 고장이 드러난다};
  % (나) two faults
  \begin{scope}[xshift=7.6cm]
    \draw[cred,line width=1.2pt,dashed] (5.2,3) -- (5.2,0.2);
    \draw (4.2,3) -- (5.2,3);
    \node[font=\scriptsize,text=cred,anchor=west,align=left] at (5.25,2.0) {첫 고장:\\아무 일 없음};
    \draw[cred,line width=1.2pt,dashed] (1.8,3) -- (1.8,0.2);
    \node[font=\scriptsize,text=cred,anchor=west,align=left] at (1.9,2.2) {둘째 고장};
    \draw[->,corange,line width=1pt] (1.95,0.45) -- (5.05,0.45);
    \node[note,text=cred,anchor=north,align=center] at (2.35,-0.45) {두 고장이 틀을 거쳐 접점을 건너뛴다\\코일이 저절로 켜진다};
  \end{scope}
\end{circuitikz}
\caption{0~V를 땅에 묶는 이유. (가) 0~V가 PE에 묶여 있으면 $+$ 쪽 선이 틀에 닿는 첫 고장이 곧 단락이어서 퓨즈가 끊긴다.
(나) 0~V가 떠 있으면 첫 고장은 드러나지 않고, 둘째 고장이 틀을 거쳐 접점을 건너뛰는 길을 만든다.}
\label{fig:p-fault}
\end{figure}

\paragraph{한 곳에서만 묶는다.} 0~V와 PE는 전원부의 연결 고리 XPE \emph{하나}로만 잇는다. 보드, 커넥터 껍데기, 모터 드라이버
어디에서도 0~V와 섀시를 직접 잇지 않는다. 두 곳에서 묶으면 금속 틀이 0~V 선과 나란한 둘째 복귀선이 된다. 코일의 복귀 전류가
전선과 틀로 나뉘어 흐르고, 그 나뉨이 케이블을 꽂고 뺄 때마다 바뀐다(\term{접지 고리}{ground loop}). 같은 이유로 26선 커넥터의
27--29번(0~V 기준선)은 공통부 쪽에서만 0~V에 묶고 스크램블러 쪽에서는 잇지 않는다. 양쪽에서 묶으면 릴레이식 유닛의 코일 복귀
전류(보드당 최대 약 0.9~A)가 메뉴 케이블로 나뉘어 흐르기 때문이다. 신호 커넥터의 0~V 핀은 신호의 \emph{기준}이지 전원의
\emph{복귀}가 아니다. 복귀는 전원 단자로 한다.

XPE는 공구로 떼어 낼 수 있는 단자다. 떼어 놓고 0~V와 PE 사이의 저항을 재면 어딘가에 둘째 접지점이 생겼는지 알 수 있다. 정상이면
보드마다의 1~M$\Omega$(아래)가 병렬로 보일 뿐이어서 12유닛에서 약 15~k$\Omega$ 이상이다(1~M$\Omega$ ÷ 약 65장, 계산). 시제품의
기준은 10~k$\Omega$ 이상이다. 0~V 쪽 선이 틀에 닿는 고장은 퓨즈를 끊지 않고 복귀 전류만 나누므로, 이 측정으로 찾는다.

\paragraph{섀시와 0~V 사이의 1~M$\Omega$ $\parallel$ 10~nF.} 보드마다 \code{CHASSIS}와 0~V 사이에 1~M$\Omega$ 저항과 10~nF(100~V)
커패시터를 나란히 한 벌 둔다. 정전기를 띤 손이 커넥터 껍데기에 닿으면 아주 빠르고 높은 펄스가 생긴다. 10~nF는 그 빠른 펄스에
짧은 길을 주어 회로의 0~V가 껍데기와 크게 벌어지지 않게 하고(\ref{sec:el-cap}절의 ``빠른 변화는 통한다''), 1~M$\Omega$은 쌓인
전하를 천천히 흘려 보낸다. 직류로는 24~V에서도 24~$\mu$A밖에 흐르지 않아(계산) 둘째 접지점이 되지 않는다.

\paragraph{차폐.} 한 함체 안의 짧은 케이블(1~m 이하)과 메뉴 케이블은 차폐하지 않는다. 받는 쪽이 수 mA, 수 ms의 릴레이 코일이어서
잡음에 둔하기 때문이다. 함체 밖으로 나가는 케이블(시험기의 케이블, 기계가 나뉘어 있을 때의 제어 케이블)은 차폐 케이블을 쓰고
차폐를 양 끝의 금속 후드(곧 섀시)에 물린다. 모터 케이블은 차폐하고 차폐를 드라이버 쪽에서 섀시에 잇는다.

\General 세탁기와 전자레인지의 전원 플러그에 접지 핀이 있는 것이 PE다. 오디오 장비를 여러 대 이었을 때 ``웅'' 하는 잡음이
나는 것은 장비마다 따로 땅에 묶여 생긴 접지 고리 때문인 경우가 많다.

\Ours 단위 시험기는 플라스틱 상자에 어댑터 전원이어서 자기 PE가 없다. 그래서 껍데기를 자기 현장 쪽 0~V에 1~M$\Omega$ $\parallel$
10~nF로 잇고, 기계에 꽂으면 케이블의 차폐를 거쳐 기계의 섀시에 닿는다. 시제품에서는 ``$+$ 쪽 선을 틀에 닿게 하면 그 갈래의
퓨즈만 끊기고 릴레이가 저절로 붙지 않는다''를 확인한다(결정 0003의 I6).

\section{비상 정지}\label{sec:p-estop}

\Why 비상 정지 회로의 핵심은 \term{폐회로 원리}{closed-circuit principle}다. 버튼이 ``정지 신호를 보내는'' 것이 아니라, 평소에
흐르는 전류를 ``끊는'' 것이다. 스위치는 모두 b접점(평소에 닫힘)이고 한 줄로 직렬이며, 그 끝에 안전 릴레이의 코일이 있다. 버튼을
누르든, 전선이 끊기든, 케이블이 빠지든, 고리의 퓨즈가 끊기든, 전원이 나가든, 결과는 같다. 고리가 열리고 안전 릴레이가 놓이고
기계가 선다. 거꾸로 버튼이 a접점으로 ``정지''를 알리는 회로라면 그 선이 끊긴 날 버튼을 눌러도 아무 일도 일어나지 않는다. 고장이
안전한 쪽으로만 나게 하는 것이다.

\General 기관사가 손을 떼면 열차가 서는 \term{데드맨 장치}{dead man's switch}가 같은 원리다. ``괜찮다''는 신호가 계속 와야 움직이고,
신호가 끊기면 이유를 묻지 않고 선다. 방화문의 자석도 평소에 전기로 문을 붙들고 있다가 전기가 끊기면 문을 놓는다.

\Ours 비상 정지 고리(\code{ESTOP})는 뱅크의 전원부가 \code{+24V\_L}에서 퓨즈 F18(0.5~A)을 거쳐 내보내고, 뱅크의 비상 정지 스위치를
지나 제어 커넥터 22번으로 나가, 공통부(패널)의 비상 정지 버튼을 지나 23번으로 돌아와 안전 릴레이 K$_\mathrm{E}$의 코일로 간다
(그림~\ref{fig:p-estop}). 제어 케이블이 빠지면 고리가 열린다. 감시 장치는 23번의 전압으로 고리가 닫혀 있는지 읽는다.

\begin{figure}[ht]
\centering
\begin{circuitikz}[x=1cm,y=1cm]
  \ctikzset{bipoles/length=8mm}
  % loop
  \node[font=\footnotesize\bfseries,anchor=west] at (-1.4,6.0) {비상 정지 고리};
  \draw (0,5.2) node[left,font=\footnotesize] {\code{+24V\_L}} to[fuse,l={F18}] (1.4,5.2) to[ncs,l={뱅크 S$_\mathrm{E1}$}] (3.0,5.2) -- (3.9,5.2)
        to[ncs,l={패널 버튼}] (5.5,5.2) -- (6.6,5.2);
  \draw[cgray,dashed,rounded corners=2pt] (3.5,4.6) rectangle (6.1,6.0);
  \node[note,anchor=north] at (4.8,4.6) {공통부(패널)};
  \node[font=\tiny,anchor=south] at (3.5,5.25) {22};
  \node[font=\tiny,anchor=south] at (6.1,5.25) {23};
  \node[rcoil,minimum width=13mm,anchor=west,fill=credbg] (KE) at (6.6,5.2) {K$_\mathrm{E}$};
  \draw (KE.east) -- (8.6,5.2) node[right,font=\footnotesize] {0 V};
  \node[note,anchor=north] at (7.25,4.75) {안전 릴레이 모듈};
  % contactor coils
  \node[font=\footnotesize\bfseries,anchor=west] at (-1.4,4.0) {접촉기 코일};
  \draw (0,3.3) node[left,font=\footnotesize] {\code{+24V\_L}} to[fuse,l={F17}] (1.4,3.3) -- (2.0,3.3) coordinate (bb)
        to[nos,l={$k_E$}] (3.6,3.3) -- (4.2,3.3);
  \node[rcoil,anchor=west] (KM1) at (4.2,3.3) {KM1};
  \draw (KM1.east) -- (6.0,3.3) node[right,font=\footnotesize] {0 V};
  \draw (bb) -- (2.0,2.3) to[nos,l_={$k_E$}] (3.6,2.3) -- (4.2,2.3);
  \node[rcoil,anchor=west] (KM2) at (4.2,2.3) {KM2};
  \draw (KM2.east) -- (6.0,2.3) node[right,font=\footnotesize] {0 V};
  \node[circ] at (bb) {};
  % main contacts
  \node[font=\footnotesize\bfseries,anchor=west] at (-1.4,1.4) {끊기는 레일};
  \draw (0,0.8) node[left,font=\footnotesize] {PSU-SW} to[nos,l={$\mathit{km}_1$}] (1.6,0.8) to[nos,l={$\mathit{km}_2$}] (3.2,0.8) -- (4.2,0.8)
        node[right,font=\footnotesize,align=left] {\code{+24V\_SW}: 계수기·시프터·순서기\\{\scriptsize (기어식: \code{+48V\_M}도 두 접촉기를 지난다)}};
  % reset
  \node[font=\footnotesize\bfseries,anchor=west] at (-1.4,-0.2) {리셋};
  \draw (0,-0.9) node[left,font=\footnotesize] {리셋 버튼} to[push button,l={S$_\mathrm{R}$}] (1.6,-0.9) to[ncs,l={$\mathit{km}_1$ (b)}] (3.2,-0.9)
        to[ncs,l={$\mathit{km}_2$ (b)}] (4.8,-0.9) -- (5.4,-0.9) node[right,font=\footnotesize] {K$_\mathrm{E}$의 리셋 입력};
\end{circuitikz}
\caption{비상 정지 사슬(전원부 문서 4.1절을 간추림). 대문자는 코일, 소문자는 그 코일이 움직이는 접점이다. 고리의 b접점 가운데
하나라도 열리면 K$_\mathrm{E}$가 놓이고, 그 접점 $k_E$가 열려 접촉기 둘이 놓이고, 직렬인 두 주 접점이 \code{+24V\_SW}를 끊는다.
다시 켜려면 리셋 버튼을 눌러야 하고, 두 접촉기가 모두 놓여 있을 때만 리셋이 된다.}
\label{fig:p-estop}
\end{figure}

부품마다 이유가 있다.
\begin{itemize}
\item \textbf{안전 릴레이 모듈.} K$_\mathrm{E}$는 보통의 릴레이가 아니라 \term{강제 안내 접점}{force-guided contacts}을 가진 안전
릴레이 모듈이다(a접점 셋, b접점 하나). 강제 안내 접점은 a접점과 b접점이 기계적으로 묶여 있어서 둘이 동시에 닫힐 수 없다. a접점이
녹아 붙으면 b접점은 열린 채 남으므로, b접점을 보면 붙은 것을 알 수 있다. 고리는 모듈의 공급 전압을 겸하고(고리 전류가 모듈의
소비 전류, 0.1~A 안팎), 시험 펄스를 내보내는 형식의 모듈은 감시 장치가 고리를 직류 전압으로 읽는 것과 맞지 않아 쓰지 않는다.
\item \textbf{접촉기를 따로 두는 이유.} 안전 릴레이 모듈의 접점은 6~A급이다. 릴레이식의 \code{+24V\_SW}(최악 23~A)나 기어식 모터의
48~V를 직접 끊을 수 없다. 그래서 모듈은 \term{접촉기}{contactor}(큰 전류용 릴레이) KM1, KM2의 코일만 움직이고, 레일은 접촉기의
주 접점이 끊는다.
\item \textbf{접촉기 둘을 직렬로.} 큰 전류를 끊는 접점은 드물게 녹아 붙는다. 하나가 붙어도 다른 하나가 끊는다. 그리고 두 접촉기의
b 보조 접점(주 접점과 기계적으로 묶인 것)이 K$_\mathrm{E}$의 리셋 회로에 직렬로 들어 있다. 어느 하나가 붙어 있으면 그 b 보조
접점이 열려 있으므로 리셋이 되지 않는다. 고장 난 채로 다시 켜지지 않는 것이다(시제품 시험 P5).
\item \textbf{수동 리셋.} 비상 정지 버튼을 풀어 놓아도 \emph{아무 일도 일어나지 않는다}. 파란 리셋 버튼을 눌러야 \code{+24V\_SW}가
돌아온다. 정전 뒤 전기가 돌아와도 마찬가지다. 버튼을 푸는 순간이나 전기가 돌아오는 순간 기계가 저절로 움직이면 안 되기
때문이다. \ref{sec:relay-latch}절의 시작/정지 회로와 같은 생각이다.
\end{itemize}

표~\ref{tab:p-events}는 사건마다 일어나는 일이다.

\begin{table}[ht]
\centering
\caption{비상 정지와 전원의 사건(전원부 문서 4.2절).}
\label{tab:p-events}
\small
\begin{tabularx}{\linewidth}{>{\raggedright\arraybackslash}p{34mm}>{\raggedright\arraybackslash}X}
\toprule
사건 & 일어나는 일 \\
\midrule
주 스위치를 켬 & \code{+24V\_L}이 산다. 레지스터, 표시, 감시 장치, 팬이 켜진다. K$_\mathrm{E}$는 놓여 있어 \code{+24V\_SW}와 \code{+48V\_M}은 없다 \\
리셋 버튼 & 고리가 닫혀 있고 두 접촉기가 놓여 있으면 K$_\mathrm{E}$와 두 접촉기가 붙는다. 릴레이식 순서기는 \code{HOME}부터 한다(계수기의 상태가 임의이므로) \\
비상 정지 버튼을 누름 & 기어식: 모터가 서고 타이밍 신호가 사라진다. 릴레이식: 모든 계수기와 시프터 코일이 놓인다. 두 기계 모두 레지스터의 유지와 램프는 남는다 \\
버튼을 풀어 놓음 & 아무 일도 없다. 리셋 버튼이 필요하다 \\
제어 케이블이 빠짐 & 비상 정지 고리와 인터록 고리가 다 열린다 \\
정전 뒤 복귀 & 리셋 버튼을 누르기 전에는 \code{+24V\_SW}가 없다 \\
\bottomrule
\end{tabularx}
\end{table}

릴레이식에는 움직이는 부위가 없어서, 비상 정지는 위험을 막는 장치라기보다 두 기계의 조작을 같게 하는 장치다. 그래도 같은 회로를
쓴다. 공통부와 단위 시험기가 뱅크의 종류를 가리지 않게 하기 위해서다.

\section{인터록: 덮개와 카트리지}\label{sec:p-interlock}

\Why \term{인터록}{interlock}은 ``이 조건이 맞을 때만 저것이 된다''를 회로로 묶는 것이다. 여기서는 ``덮개가 닫혀 있을 때만 돈다''다.
구조는 비상 정지와 같다. 덮개 스위치(b접점)의 직렬 고리가 안전 릴레이 K$_\mathrm{I}$를 먹이고, 고리가 열리면 K$_\mathrm{I}$가 놓인다.
다른 점은 무엇을 끊느냐와 다시 닫혔을 때 어떻게 되느냐다(표~\ref{tab:p-loops}).

\General 전자레인지는 문을 열면 마그네트론의 전원이 문 스위치의 접점으로 끊긴다. 세탁기는 도는 동안 문이 잠긴다. 둘 다 사람이 위험한
곳에 손을 넣을 수 있는 순간에 위험을 없앤다.

\Ours 두 기계의 인터록은 막는 것이 다르다.

\paragraph{기어식: 덮개를 열면 모터의 48~V만.} 고리는 스위치 셋의 직렬이다. 앞 덮개 양 끝의 덮개 스위치 S$_\mathrm{I1}$, S$_\mathrm{I2}$와
손 크랭크 구멍의 셔터 스위치 S$_\mathrm{I3}$이다. 고리가 열리면 K$_\mathrm{I}$가 놓이고 접촉기 KM3이 놓여 \code{+48V\_M}만 끊는다
(그림~\ref{fig:p-unit}). 시험 전압과 레지스터는 살아 있다. 덮개를 닫으면 K$_\mathrm{I}$는 저절로 다시 붙는다(리셋 회로가 KM3의 b
보조 접점뿐인 자동 리셋). 그러면 덮개를 닫는 순간 드럼이 다시 돌기 시작하지 않을까. 이를 막는 것이 \code{V48\_OK}다.
\begin{itemize}
\item KM3에는 a 보조 접점(13--14)이 하나 더 있어, KM3이 붙어 있는 동안 \code{+24V\_L}을 \code{V48\_OK}로 구동부 제어 보드에 보낸다
(2.2~k$\Omega$ + 광커플러, 약 10~mA).
\item 제어 보드는 \code{V48\_OK}가 꺼졌다 켜지면 \code{RUN}이 한 번 내려갔다 올라와야(시작 스위치를 껐다 켜야) 움직인다. 이 규칙은 제어
보드의 펌웨어 모델에 있다(\code{RUNARM}, \ref{sec:gk-safety}절).
\item a 보조 접점은 안전 기능이 아니다. 48~V를 끊는 것은 KM3의 주 접점이다. 다만 a 접점이 붙어 버리면 ``다시 켜야 움직인다''는 규칙이
사라지므로, a와 b가 기계적으로 묶인 보조 접점 블록을 고른다(IEC 60947-5-1 부속서 L, 원문 미확인). 그러면 a가 붙어 있는 동안 b가
닫히지 않아 K$_\mathrm{I}$가 리셋되지 않으므로 고장이 드러난다.
\end{itemize}

\paragraph{기어식: 덮개를 열면 빠른 축을 잡는 브레이크.} 덮개가 열려 모터에 힘이 없는 동안 관람객이 빠른 줄의 드럼을 손으로
돌리면 자리올림이 가운데 줄을 넘겨 크립 오프셋이 \emph{조용히} 틀릴 수 있다. 그래서 F 축에 무여자 작동 브레이크를 단다. 그 까닭과
정지 토크(4~N$\cdot$m 이상, 억지로 비틀면 축이 아니라 드럼의 인덱스가 넘어간다)는 \ref{sec:gm-brake}절에서 보았다(그림~\ref{fig:p-interlock}의
(나)). 전원과 안전 쪽에서 덧붙일 것은 이렇다.
\begin{itemize}
\item 전기가 없으면 잡으므로 폐회로 원리와 같은 방향이다.
\item 코일은 \code{+24V\_L}에서 F22(1~A)를 거쳐 덮개 스위치 S$_\mathrm{I1}$, S$_\mathrm{I2}$의 \emph{둘째} b접점을 직렬로 지나 먹는다.
덮개가 닫혀 있으면 풀리고, 열리면 잡는다. 크랭크 셔터는 넣지 않아 덮개를 닫은 채의 손 크랭크 시연은 그대로 된다.
\item 운전 중에 덮개를 열면 브레이크가 F를 약 9~ms 안에 세운다(계산: 관성 약 5.5×10$^{-3}$~kg$\cdot$m$^2$, 6.3~rad/s, 4~N$\cdot$m에서
8.6~ms). 브레이크가 없던 설계에서는 48~V만 끊겨 관성으로 섰다.
\item 안전 기능은 바뀌지 않는다. 48~V를 끊는 것은 여전히 K$_\mathrm{I}$와 KM3이다. 브레이크는 덮개 스위치의 둘째 접점에 매달린 보조
기능이어서, 브레이크가 풀린 채로 고장 나도 브레이크가 없던 설계와 같은 상태가 될 뿐이다.
\end{itemize}

\paragraph{릴레이식: 카트리지 덮개를 열면 순서기의 전원.} 릴레이식에는 움직이는 위험 부위가 없어 원래 인터록 고리를 뱅크 안에서
점퍼로 닫아 두었다. 관람객이 로터 카트리지를 바꾸게 되면서(\code{rotor-cartridge.md}), 이 고리가 카트리지의 \emph{검출 고리}가 되었다.
받침 기판의 자리마다 검출 패드 한 쌍(카트리지 안에서 서로 이어진 패드에 닿는 스프링 핀)과 누름 덮개의 스위치를 직렬로 잇고, 받침
기판들을 유닛 순서로 이어 한 고리로 만든다. 카트리지가 바로 앉고 덮개가 잠겨야만 닫힌다. 고리가 열리면 K$_\mathrm{I}$가 놓이고, 그
a접점이 클럭 순서기의 전원 갈래(F38, 최대 약 1.7~A)를 직접 끊는다. 접촉기는 두지 않는다(모듈 접점의 직류 차단 정격은 미확인).
\begin{itemize}
\item 순서기가 서면 \code{STROBE}, \code{JUDGE}, \code{CLEAR}, \code{CK}/\code{NCK}가 모두 사라진다. 시험 전압은 레지스터의
\code{STROBE} 접점을 거쳐서만 26선에 가므로, \emph{모든 스크램블러 선에 전압이 없다}. 빈 자리의 스프링 핀에 손이 닿아도 전압이 없다.
\code{CK} = \code{NCK} = 0이므로 계수기와 시프터 코일도 놓인다.
\item \textbf{왜 \code{RUN}이 아니라 전원을 끊는가.} 한 걸음 버튼은 \code{RUN}과 무관하게 한 위치를 시험한다(\ref{sec:t-sequencer}절).
\code{RUN}만 막았다면 덮개가 열린 채 한 걸음 버튼으로 \code{STROBE}를 낼 길이 남는다. 한 걸음 버튼의 공급도 순서기 전원에서 나오므로,
전원을 끊으면 그 길도 함께 사라진다.
\item 다시 잠그면 K$_\mathrm{I}$가 저절로 붙지만, 순서기는 \code{HOME}을 하기 전에는 \code{RUN}을 무시한다. 덮개를 닫는 것만으로 다시
돌지 않는다. 기어식의 \code{V48\_OK} 규칙과 같은 뜻이다.
\item \textbf{먼저 끊고 나중에 뗀다.} 걸쇠를 풀기 시작하는 첫 10° 안에서 스위치가 열리고, 패드가 핀에서 떨어지는 것은 약 45° 뒤다. 손으로
푸는 속도에서 그 사이가 약 0.1초이고, 순서기 전원은 약 30~ms 안에 끊긴다(모두 추정). 꽂을 때는 반대로 패드가 먼저 닿고 스위치가
마지막에 닫힌다. \ref{sec:relay-bbm}절의 ``먼저 끊고 나중에 잇기''를 손의 동작으로 만든 것이다.
\item 이 고리는 \emph{안전 기능이 아니라 기능 인터록}이다. 24~V이고 움직이는 것이 없으므로 위험은 없다. 막는 것은 반쯤 꽂힌 카트리지를
지나는 \code{STROBE}(거짓 정지, 핀과 패드 사이의 작은 불꽃), 운전 중에 로터를 바꿔 결과가 섞이는 것, 빈 자리의 핀에 전압이 있는 것이다.
그래도 배선, 감시, 리셋을 기어식과 같게 두어 공통부와 시험기가 뱅크를 가리지 않는다.
\end{itemize}

\begin{figure}[tbp]
\centering
\resizebox{\linewidth}{!}{% 로터를 바꾸는 동안의 인터록: 릴레이식 카트리지 검출 고리(가), 기어식 브레이크(나). 그림 fig:hw-rotor-interlock
% 근거: hardware/common/rotor-cartridge.md 3.8절, 4.3절, hardware/common/power-unit.md 4.1절, docs/decisions/0003-interfaces.md 4.6절.
\begin{tikzpicture}[x=1cm,y=1cm,
    el/.style={draw,rounded corners=1pt,minimum height=6.5mm,minimum width=8mm,inner sep=1.5pt,font=\tiny,align=center,fill=white},
    det/.style={el,draw=corange,fill=corangebg},
    lid/.style={el,draw=cblue,fill=cbluebg},
    fu/.style={el,draw=cgray,fill=clight},
    coil/.style={el,draw=black,fill=cyellowbg},
    kc/.style={el,draw=cred,fill=credbg},
    ln/.style={line width=0.6pt},
    lab/.style={font=\tiny,align=left,inner sep=1pt}]
  % ================= (가) 릴레이식 =================
  \node[font=\footnotesize\bfseries,anchor=west] at (-0.7,0.95) {(가) 릴레이식: 카트리지의 검출 고리가 클럭 순서기의 전원을 끊는다};
  \node[lab,anchor=east] (v1) at (0.45,0) {\code{+24V\_L}};
  \node[fu] (f19) at (1.15,0) {F19\\0.5 A};
  \node[det] (dl) at (2.4,0) {검출\\L};
  \node[lid] (sl) at (3.4,0) {덮개\\L};
  \node[det] (dm) at (4.4,0) {검출\\M};
  \node[lid] (sm) at (5.4,0) {덮개\\M};
  \node[det] (dr) at (6.4,0) {검출\\R};
  \node[lid] (sr) at (7.4,0) {덮개\\R};
  \draw[cgray,dashed,rounded corners=2pt] (1.9,-0.5) rectangle (7.9,0.5);
  \node[lab,anchor=south west,text=cgray] at (1.9,0.5) {유닛 7의 받침 기판 (J16 입출력, J13--J15 덮개 스위치)};
  \node[font=\scriptsize] (dots) at (8.3,0) {$\cdots$};
  \node[el] (u12) at (9.1,0) {유닛\\12};
  \node[el] (mt) at (10.3,0) {정비 덮개\\(점퍼)};
  \node[el] (cc) at (12.2,0) {제어 커넥터 24\\$\Rightarrow$ 공통부 $\Rightarrow$ 25};
  \node[coil] (ki) at (13.95,0) {K$_\mathrm{I}$\\코일};
  \node[lab,anchor=west] (z1) at (14.55,0) {0 V};
  \draw[ln] (v1) -- (f19) -- (dl) -- (sl) -- (dm) -- (sm) -- (dr) -- (sr) -- (dots) -- (u12) -- (mt) -- (cc) -- (ki) -- (z1);
  % 순서기의 전원 갈래
  \node[lab,anchor=east] (v2) at (0.45,-1.75) {PSU-SW};
  \node[kc] (km) at (1.3,-1.75) {KM1·KM2};
  \node[fu] (f38) at (2.65,-1.75) {F38\\3.15 A};
  \node[el,draw=cred] (kic) at (3.9,-1.75) {K$_\mathrm{I}$\\a접점};
  \node[el,draw=cblue,fill=cbluebg,minimum width=17mm] (seq) at (5.55,-1.75) {클럭 순서기};
  \draw[ln] (v2) -- (km) -- (f38) -- (kic) -- (seq);
  \node[lab,anchor=west] (out) at (7.0,-1.75) {\code{STROBE}, \code{JUDGE}, \code{CLEAR}, \code{CK}, \code{NCK}와\\한 걸음 버튼의 공급(패널 입력)이 함께 사라진다};
  \draw[->,ln] (seq) -- (out);
  \draw[->,cred,dashed,line width=0.6pt] (ki.south) -- (13.95,-0.95) -- (3.9,-0.95) -- (kic.north);
  \node[lab,text=cred,anchor=south] at (10.6,-0.93) {고리가 열리면 K$_\mathrm{I}$가 놓인다 (다시 닫히면 스스로 붙지만 순서기는 원위치를 기다린다)};

  % ================= (나) 기어식 =================
  \node[font=\footnotesize\bfseries,anchor=west] at (-0.7,-2.85) {(나) 기어식: 덮개 스위치의 둘째 b접점이 빠른 라인 축 F의 브레이크를 놓는다};
  \node[lab,anchor=east] (v3) at (0.45,-3.6) {\code{+24V\_L}};
  \node[fu] (f19b) at (1.15,-3.6) {F19\\0.5 A};
  \node[lid] (s1a) at (2.6,-3.6) {S$_\mathrm{I1}$ 덮개\\첫째 b};
  \node[lid] (s2a) at (4.1,-3.6) {S$_\mathrm{I2}$ 덮개\\첫째 b};
  \node[el] (s3) at (5.6,-3.6) {S$_\mathrm{I3}$\\크랭크 셔터};
  \node[el] (cc2) at (7.3,-3.6) {제어 커넥터\\24 $\Rightarrow$ 25};
  \node[coil] (ki2) at (8.75,-3.6) {K$_\mathrm{I}$\\코일};
  \draw[ln] (v3) -- (f19b) -- (s1a) -- (s2a) -- (s3) -- (cc2) -- (ki2);
  \node[lab,anchor=west] at (9.35,-3.6) {0 V. K$_\mathrm{I}$ $\to$ KM3: 모터의 48 V만 끊음(그대로)};
  \node[lab,anchor=east] (v4) at (0.45,-4.85) {\code{+24V\_L}};
  \node[fu] (f22) at (1.15,-4.85) {F22\\1 A};
  \node[lid] (s1b) at (2.6,-4.85) {S$_\mathrm{I1}$ 덮개\\둘째 b};
  \node[lid] (s2b) at (4.1,-4.85) {S$_\mathrm{I2}$ 덮개\\둘째 b};
  \node[coil,minimum width=14mm] (brk) at (5.75,-4.85) {브레이크 코일\\(24 V)};
  \node[lab,anchor=west] (z4) at (6.6,-4.85) {0 V};
  \draw[ln] (v4) -- (f22) -- (s1b) -- (s2b) -- (brk) -- (z4);
  \draw[cgray,dashed,line width=0.4pt] (s1a.south) -- (s1b.north);
  \draw[cgray,dashed,line width=0.4pt] (s2a.south) -- (s2b.north);
  \node[lab,text=cgray,anchor=west] at (4.6,-4.22) {같은 스위치의 두 접점};
  \node[lab,anchor=west] at (7.25,-4.85) {덮개가 닫혀 있으면 코일에 전류가 흘러 브레이크가 풀린다.\\덮개를 열면 스프링이 F를 잡는다(4 N·m 이상).\\크랭크 셔터는 넣지 않는다: 덮개를 닫은 채 손 크랭크 시연};
  % 브레이크가 없을 때 상태가 바뀌는 길
  \node[el,minimum width=20mm] (drm) at (1.0,-6.45) {빠른 줄 드럼\\(관람객이 돌림)};
  \node[el,draw=cgray,fill=clight,minimum width=16mm] (F) at (3.6,-6.45) {라인 축 F};
  \node[el,draw=cgray,fill=clight,minimum width=16mm] (C) at (6.2,-6.45) {주기 축 C};
  \node[el,draw=cgray,fill=clight,minimum width=18mm] (M) at (10.2,-6.45) {가운데 줄 M};
  \draw[->,ln] (drm) -- node[above,lab] {마이터} (F);
  \draw[->,ln] (F) -- node[above,lab] {26:30} (C);
  \draw[->,ln] (C) -- node[above=2.5mm,lab,align=center] {자리올림 1\\(F 1.15바퀴에 한 칸)} (M);
  \node[kc] (brk2) at (3.6,-7.35) {브레이크};
  \draw[cred,line width=0.8pt] (brk2.north) -- (F.south);
  \node[lab,anchor=west] at (11.2,-6.45) {브레이크가 없으면 M이 넘어가\\나중에 끼운 드럼의 글자가\\원위치와 어긋난다};
  \node[lab,anchor=west,text=cred] at (4.25,-7.35) {잡으면 F, C가 서고 M, S는 잠금 원호가 잡는다};
\end{tikzpicture}
}
\caption{로터를 바꾸는 동안의 인터록(보고서와 같은 그림, \code{rotor-cartridge.md} 3.8·4.3절). (가) 릴레이식: 카트리지의 검출 고리가
K$_\mathrm{I}$를 거쳐 클럭 순서기의 전원을 끊는다. (나) 기어식: 덮개 스위치의 첫째 b접점은 인터록 고리에, 둘째 b접점은 F 축 브레이크의
코일에 들어간다. 맨 아래 줄은 브레이크가 없을 때 드럼을 돌리면 기계 상태가 바뀌는 길이다.}
\label{fig:p-interlock}
\end{figure}

\begin{table}[ht]
\centering
\caption{세 고리의 비교(결정 0003 4.6절, 전원부 문서 4절, \code{rotor-cartridge.md} 3.8절).}
\label{tab:p-loops}
\small
\begin{tabularx}{\linewidth}{>{\raggedright\arraybackslash}p{19mm}>{\raggedright\arraybackslash}X>{\raggedright\arraybackslash}X>{\raggedright\arraybackslash}X}
\toprule
& 비상 정지 (두 기계) & 덮개 인터록 (기어식) & 카트리지 인터록 (릴레이식) \\
\midrule
고리의 스위치 & 뱅크의 버튼 + 패널의 버튼 & 앞 덮개 2 + 크랭크 셔터 & 자리마다 검출 패드 + 덮개 스위치, 유닛 순서로 \\
받는 것 & K$_\mathrm{E}$ & K$_\mathrm{I}$ & K$_\mathrm{I}$ \\
끊는 것 & KM1·KM2 $\to$ \code{+24V\_SW}(기어식은 \code{+48V\_M}도) & KM3 $\to$ \code{+48V\_M}만 & K$_\mathrm{I}$ a접점 $\to$ 순서기 갈래 \\
남는 것 & \code{+24V\_L}: 레지스터의 유지와 램프, 표시, 감시 & 시험 전압과 레지스터(손 크랭크 시연) & 레지스터의 표시 \\
다시 닫으면 & 아무 일 없음. 리셋 버튼 & K$_\mathrm{I}$ 자동 복귀. \code{RUN}을 껐다 켜야 진행 & K$_\mathrm{I}$ 자동 복귀. \code{HOME} 전에는 진행 없음 \\
성격 & 안전 기능 & 안전 기능 & 기능 인터록 \\
\bottomrule
\end{tabularx}
\end{table}

\section{왜 하드웨어로 하는가}\label{sec:p-hardware}

\Why 덮개가 열렸는지는 마이크로컨트롤러도 읽을 수 있다. 그런데 왜 접점의 고리와 접촉기로 하는가. 프로그램은 멈출 수 있고(무한
루프, 전원 잡음으로 인한 재시작), 잘못 고친 판이 올라갈 수 있고, 드물게 지나는 갈래에 버그가 숨어 있을 수 있다. 그리고 그 실패는
대개 ``아무것도 하지 않음''이어서, 전원을 끊어야 할 때 끊지 않는다. 직렬 접점의 고리는 상태가 없다. 고장 나는 방식이 몇 가지뿐이고
(전선이 끊김, 접점이 안 닫힘, 전원이 나감), 그 대부분이 ``고리가 열림'', 곧 정지다. 남는 위험한 고장(접점이 녹아 붙음)은 직렬 접촉기와
리셋 회로의 b 보조 접점이 드러낸다. 안전 기능의 몸은 ``에너지를 없앤다''이고, 그것은 접점이 물리적으로 떨어지는 것으로 한다.

\General 전자레인지의 문 스위치는 프로그램이 아니라 마그네트론의 전원선에 직렬로 들어 있다. 조작판의 프로그램이 멈춰도 문을 열면
꺼진다.

\Ours 이 프로젝트에서 소프트웨어는 안전 고리의 \emph{옆}에 있다. 무엇이 어느 쪽인지 적어 보면 이렇다.
\begin{itemize}
\item \textbf{하드웨어.} 비상 정지와 인터록의 고리, 안전 릴레이 모듈, 접촉기, 퓨즈, 레일의 분리, 브레이크, 카트리지의 검출 고리. 그리고
판정 자체(\ref{sec:t-register}절의 테스트 레지스터)와 기어식의 타이밍 경로(광센서에서 커넥터 핀까지 펌웨어가 없다, \ref{sec:t-gear}절).
\item \textbf{소프트웨어.} 감시 장치는 23번과 25번 핀의 전압으로 고리를 읽어 ``카트리지 덮개를 잠그세요'' 같은 안내를 띄운다. 알릴 뿐 고리를
닫지 않는다. 설계 개요는 감시 장치를 ``판정·구동 논리에는 끼지 않는다(모터 드라이버 제어는 예외)''고 적는다. 기어식 제어 보드의 펌웨어는 모터의 가감속과 되돌림을
맡고, 덮개를 닫은 뒤 \code{RUN}을 다시 올려야 움직이는 규칙(\code{RUNARM})도 펌웨어 모델에 있다.
\end{itemize}
마지막 줄이 경계를 보여 준다. 덮개를 닫은 뒤 저절로 돌지 않게 하는 규칙은 펌웨어에 있지만, 그 규칙이 틀려도 ``덮개가 열려 있는 동안
48~V가 없다''는 사실은 바뀌지 않는다. 사람의 손이 닿을 수 있는 동안의 보호는 하드웨어만으로 성립하고, 소프트웨어는 그 위에서 편의와
안내를 맡는다.

\section{정직하게: 단일 채널, ISO 13849-1 미평가}\label{sec:p-honest}

\Why 기계의 안전 회로가 ``얼마나 믿을 만해야 하는가''를 정하는 국제 규격이 ISO 13849-1이다~\cite{iso13849-1}. 대략의 흐름은 이렇다. 먼저
위험을 평가해(다칠 수 있는 정도, 위험에 노출되는 빈도, 피할 수 있는 가능성) 안전 기능마다 요구되는 \term{성능 수준}{performance level,
PL}을 정한다. 그리고 회로의 구조(몇 채널인가, 고장을 스스로 찾아내는가)와 부품의 신뢰도로 그 수준에 이르는지 보인다. 높은 수준일수록
보통 \term{이중 채널}{dual channel}을 요구한다. 같은 고리를 두 벌 두고 안전 릴레이가 두 채널이 서로 같은지 늘 비교하는 것이다.

\term{단일 채널}{single channel}이 무엇을 놓치는지 예를 들면 이렇다. 고리의 두 선이 케이블 안에서 피복이 벗겨져 서로 붙으면, 그 뒤로는
비상 정지 버튼을 눌러도 고리가 열리지 않는다. 고장 하나가 안전 기능을 조용히 없앤 것이다. 이중 채널이라면 한 채널만 단락되었을 때 두
채널이 달라지므로 모듈이 알아채고 리셋을 거부한다.

\General 자동차의 브레이크 배관이 앞뒤 두 계통으로 나뉜 것이 이중 채널의 예다. 한 계통이 새도 다른 계통으로 선다.

\Ours 이 프로젝트의 안전 회로는 \emph{단일 채널}이고, ISO 13849-1의 요구 성능 수준 판정은 \emph{하지 못했다}. 전원부 문서와 참고 문헌이
그렇게 적고 있고, 이 문서도 그 이상을 말하지 않는다.
\begin{itemize}
\item \textbf{왜 단일 채널인가.} 제어 커넥터에 고리마다 핀이 둘뿐이다(결정 0003 5.2절). 이 기계에 남는 위험은 작다고 보았다. 기어식 드럼의
둘레 힘은 약 20~N으로 묶여 있고(결정 0002), 릴레이식에는 움직이는 부위가 없다.
\item \textbf{무엇을 하지 않았는가.} 규격 원문을 열람하지 못했고, 위험 평가와 성능 수준 판정을 하지 않았다. 판정에서 이중 채널이 필요하다는
결론이 나오면, 제어 커넥터의 예약 핀(37번, 그리고 기계에서는 쓰지 않는 26--27번)을 둘째 채널로 쓰는 것을 검토한다.
\item \textbf{원문으로 확인하지 못한 것들.} 전원 장치의 SELV 표시와 48~V가 그 범위인지, 제어 회로 접지(IEC 60204-1), 기계적으로 묶인 보조
접점(IEC 60947-5-1 부속서 L), 접촉기의 직류 정격, K$_\mathrm{I}$ 접점의 직류 차단 정격, 브레이크와 스프링 핀의 품번. 문서마다
``미확인''으로 적혀 있다.
\item \textbf{전시 장소가 정해지면.} 그곳의 전기·기계 안전 기준으로 다시 본다. 전원 장치의 인증(국내면 KC)도 구매 때 확인한다.
\end{itemize}
이 장이 설명한 것은 \emph{설계의 의도}와 \emph{그 의도의 이유}다. 인증받은 안전이 아니다. 판정 결과에 따라 둘째 채널을 낼 핀은 남겨 두었다.

\section{시제품에서 확인할 것}\label{sec:p-tests}

전원부는 회로도 이전이므로, 이 장의 주장 대부분은 시제품에서 잰다(표~\ref{tab:p-tests}).

\begin{table}[ht]
\centering
\caption{전원부의 확인 항목(전원부 문서 8절).}
\label{tab:p-tests}
\small
\begin{tabularx}{\linewidth}{l>{\raggedright\arraybackslash}X>{\raggedright\arraybackslash}X}
\toprule
& 항목 & 기준 \\
\midrule
P1 & 최대 부하에서 보드 단자의 전압 & 23.0~V 이상 \\
P2 & $+$ 쪽 선을 틀에 닿게 함 & 그 갈래의 퓨즈만 끊기고 릴레이가 저절로 붙지 않음 \\
P3 & XPE를 뗀 상태의 0~V--PE 절연 저항 & 10~k$\Omega$ 이상 \\
P4 & 비상 정지에서 \code{+24V\_SW}와 \code{+48V\_M}이 사라지는 시간 & 0.1초 이하(기어식) \\
P5 & 접촉기 하나의 주 접점을 붙여 놓고 리셋 & 리셋되지 않음 \\
P6 & 비상 정지를 풀어 놓았을 때, 덮개를 닫았을 때 & 저절로 움직이지 않음 \\
P7 & 릴레이식에서 다이얼을 모두 같은 값(15 또는 23)으로 두고 구동 & 전원 장치가 전류 제한에 들지 않고 유닛의 퓨즈가 끊기지 않음 \\
P8 & 전원 장치 세 대의 돌입 전류 & 차단기가 떨어지지 않음 \\
\bottomrule
\end{tabularx}
\end{table}

P5와 P2가 이 장의 두 핵심을 시험한다. P5는 ``고장 난 채로 다시 켜지지 않는다''를, P2는 ``고장이 숨지 않고 드러난다''를 본다. 음성 시험이
검증기를 일부러 틀린 회로로 시험하듯(\ref{sec:v-mutants}절), 안전 회로도 일부러 고장을 내서 시험한다.

\section{정리}

\begin{itemize}
\item 220~V는 인증된 전원 장치 안에만 둔다. 관람객이 닿는 것은 모두 24~V 직류이고, 48~V는 기어식 구동부의 덮개 안에만 있다.
\item 레일은 끊어야 하는 때로 나눈다. \code{+24V\_L}은 결과를 붙들고 보이는 쪽이라 주 스위치만 끊고, \code{+24V\_SW}는 비상 정지가, \code{+48V\_M}은
비상 정지와 덮개가 끊는다.
\item 전원 장치는 평균이 아니라 최악으로 고른다. 다이얼 설정 하나로 최악이 실제로 생기기 때문이다. 아직 모르는 코일 전력이 크기를 바꿀 수 있다.
\item 보호는 두 겹이다. 보드 안의 고장은 보드의 자동 복귀 퓨즈가, 전선의 고장은 전원부의 퓨즈가 맡는다. 보드 입구에는 자동 복귀 퓨즈,
P채널 MOSFET 역극성 보호, TVS, 커패시터가 같은 순서로 있다.
\item 0~V는 PE에 묶되 한 곳에서만 묶는다. 묶어야 $+$ 쪽 고장이 퓨즈를 끊어 드러나고, 한 곳이어야 접지 고리가 없다. 섀시는 0~V가 아니며
1~M$\Omega$ $\parallel$ 10~nF로만 만난다.
\item 비상 정지와 인터록은 b접점의 직렬 고리다(폐회로 원리). 끊기는 모든 고장이 정지가 된다. 안전 릴레이 모듈이 직렬 접촉기 둘을 움직이고,
붙은 접촉기는 리셋을 막는다. 무엇도 저절로 다시 돌지 않는다.
\item 기어식의 덮개는 48~V를 끊고 브레이크로 빠른 축을 잡는다. 릴레이식의 카트리지 고리는 순서기의 전원을 끊어 시험 전압과 한 걸음 버튼을
함께 없앤다. 앞의 것은 안전 기능, 뒤의 것은 기능 인터록이다.
\item 사람을 지키는 일은 하드웨어가 하고 소프트웨어는 옆에서 알린다. 안전 회로는 단일 채널이고 ISO 13849-1 평가는 하지 않았다. 이 장은
설계의 의도와 이유이지 인증이 아니다.
\end{itemize}
